\section*{Capítulo \ref{ch:g9:periodic-table-first-look} --- A tabela periódica: um primeiro olhar}

\begin{solution}{exo:g9:periodic-table-first-look:1}
Em ordem crescente de \omterm{def:g9:inside-the-atom:atomic-number}{número atômico} $Z$.
\end{solution}

\begin{solution}{exo:g9:periodic-table-first-look:2}
Carbono: \omterm{def:g9:periodic-table-first-look:period}{período} 2, \omterm{def:g9:periodic-table-first-look:period}{grupo} 14. Magnésio: \omterm{def:g9:periodic-table-first-look:period}{período} 3, \omterm{def:g9:periodic-table-first-look:period}{grupo} 2. Argônio:
\omterm{def:g9:periodic-table-first-look:period}{período} 3, \omterm{def:g9:periodic-table-first-look:period}{grupo} 18.
\end{solution}

\begin{solution}{exo:g9:periodic-table-first-look:3}
\omterm{def:g9:periodic-table-first-look:metal}{Metais}: ferro, cobre, alumínio. \omterm{def:g9:periodic-table-first-look:metal}{Não metais}: enxofre, oxigênio, cloro.
\end{solution}

\begin{solution}{exo:g9:periodic-table-first-look:4}
O potássio (\ce{K}) abaixo do sódio; o bromo (\ce{Br}) abaixo do cloro.
\end{solution}

\begin{solution}{exo:g9:periodic-table-first-look:5}
O cálcio (\ce{Ca}, $Z = 20$) e o nitrogênio (\ce{N}, $Z = 7$).
\end{solution}

\begin{solution}{exo:g9:periodic-table-first-look:6}
Ele está no mesmo grupo que o sódio, o \omterm{def:g9:periodic-table-first-look:period}{grupo} 1, e os elementos de um
grupo se comportam igual.
\end{solution}

\begin{solution}{exo:g9:periodic-table-first-look:7}
Um \omterm{def:g9:periodic-table-first-look:metal}{metal}: ele é brilhante, conduz a eletricidade e forma um \omterm{def:g9:ions:ion}{íon} positivo.
\end{solution}

\begin{solution}{exo:g9:periodic-table-first-look:8}
Para manter os elementos de propriedades parecidas na mesma coluna,
mesmo que isso significasse deixar uma casa vazia. As lacunas eram
previsões: os elementos que faltavam foram encontrados mais tarde com as
propriedades que ele tinha descrito.
\end{solution}

\begin{solution}{exo:g9:periodic-table-first-look:9}
Carbono: $\ce{C} = 12$; oxigênio: $\ce{O} = 16$. As lacunas ``? $= 68$'' (ao lado de
Al $= 27.4$) e ``? $= 70$'' (ao lado de Si $= 28$).
\end{solution}

\begin{solution}{exo:g9:periodic-table-first-look:10}
\omterm{def:g9:periodic-table-first-look:period}{Período} 2: 8 elementos (do lítio ao neônio). \omterm{def:g9:periodic-table-first-look:period}{Período} 4: 18 elementos (do
potássio ao criptônio).
\end{solution}

\begin{solution}{exo:g9:periodic-table-first-look:11}
\omterm{def:g9:periodic-table-first-look:period}{Grupo} 18. Onde nada pode reagir: o argônio dentro das lâmpadas e no gás
de soldagem, o hélio nos balões (ele não \omterm{def:g4:what-a-fire-needs:burning}{queima}).
\end{solution}

\begin{solution}{exo:g9:periodic-table-first-look:12}
$(27.0 + 114.8)/2 = 70.9$, perto de 69.7.
\end{solution}

\begin{solution}{pb:g9:periodic-table-first-look:1}
\textbf{1.} $Z = 32$; \omterm{def:g9:periodic-table-first-look:period}{período} 4, \omterm{def:g9:periodic-table-first-look:period}{grupo} 14.

\textbf{2.} Em cima: o silício. Embaixo: o estanho. À esquerda: o gálio.
À direita: o arsênio.

\textbf{3.} Um metaloide.

\textbf{4.} Ele está no mesmo grupo que o silício, e os elementos de um
grupo se comportam igual.

\textbf{5.} $(28.1 + 118.7)/2 = 73.4$.

\textbf{6.} $(69.7 + 74.9)/2 = 72.3$.

\textbf{7.} A média esquerda--direita (72.3): ao longo de um \omterm{def:g9:periodic-table-first-look:period}{período}, o
peso atômico cresce de maneira regular, elemento por elemento, enquanto
ao descer num grupo ele dá grandes saltos.

\textbf{8.} $72.6 - 70 = 2.6$.

\textbf{9.} $(28.1 + 118.7 + 69.7 + 74.9)/4 = 291.4/4 = 72.85$, isto é,
72.9.

\textbf{10.} A média, 72.9, está perto ao mesmo tempo dos 72 de
Mendeleiev e dos 72.3 de Winkler (e dos 72.6 de hoje).

\textbf{11.} Um elemento previsto antes que alguém o tivesse visto foi
encontrado com as propriedades anunciadas: a tabela não era só uma
lista, ela podia fazer previsões.

\textbf{12.} Cerca de 72.9.
\end{solution}
